\documentclass[10pt,a4paper]{amsart}
\usepackage[margin=19mm]{geometry}
\usepackage{amsmath,amssymb,mathtools}
\usepackage[scr=rsfs,cal=euler]{mathalfa}
\usepackage{xfrac}
\usepackage[unicode,hidelinks,bookmarks=false]{hyperref}
\newcommand{\ie}{i.e.\ }
\newcommand{\cf}{cf.\ }
\newcommand{\resp}{respectively\ }
\newcommand{\Subsection}{\subsection}
\providecommand{\nobreakdash}{\nobreak\hbox{-}}
\setlength{\emergencystretch}{2em}
\multlinegap0pt
\usepackage{amssymb}
\usepackage[shortlabels]{enumitem}
\usepackage{bm}
\newtheorem{lem}{Lemma}
\newcommand{\C}{\mathbb{C}}
\newcommand{\N}{\mathbb{N}}
\newcommand{\cF}{\mathcal{F}}
\newcommand{\fm}{\mathfrak{m}}
\DeclareMathOperator{\ord}{ord}
\title{On Positivity of the Limit $F$-signature}
\author{Yuchen Liu, Suchitra Pande}
\date{Source: arXiv:2605.16636v1 (CC BY 4.0)}
\begin{document}
\maketitle

\noindent\emph{Source and licence.} On Positivity of the Limit $F$-signature. Version arXiv:2605.16636v1 (CC BY 4.0), distributed by arXiv under the Creative Commons Attribution 4.0 International licence, http://creativecommons.org/licenses/by/4.0/, as recorded in the arXiv licence metadata for this exact version and shown on the abstract page https://arxiv.org/abs/2605.16636v1. Full licence notice: \texttt{licenses/arxiv-2605.16636-CC-BY-4.0.txt}.\par\smallskip

\noindent\textit{Editorial scope.} Counted unit: the article's lem fromCtocharp (source0.tex lines 1334--1336). The article's complete original proof of it is delivered with it (source0.tex lines 1338--1363, one \texttt{proof} environment of 26 lines). No mathematics is written, repaired or re-derived: the statement and the proof are the article's own lines, in the article's own notation, and no retained line is substituted or re-flowed.  No further source-local context is needed: every cross-reference of every retained line resolves inside the two delivered spans.  Nothing was omitted from any retained span, so the package carries no omission record.  No retained line contains a citation, so the wrapper carries no bibliography and no bibliography entry.  No retained line contains a figure environment, an image-inclusion command or drawing code, so no image is delivered and the package declares no asset dependency.  Editorial formatting only, in addition: an AMS \texttt{amsart} preamble that restates the article's own declarations and macros used by the retained lines, namely the environments \texttt{lem} on separate counters, together with the standard packages those lines need.  The retained environments print in this document as Lemma 1 in order of appearance, whereas the article prints the same items with the numbering of its own full document.  The complete native source of the article as distributed in the arXiv e-print for this exact version is bundled unchanged as \texttt{sources/200732/source0.tex}.
\begin{lem} \label{fromCtocharp}
    For any $m'$ divisible enough, suppose that there exists a section $g_\C \in S_{\C,m m'}$ such that $g_\C$ is divisible by $f_{\C}^{ \alpha_0 m' }$ and $\ord_{E_\C} (g_\C) = N$. Then, there exists $g \in S_{m m'}$ such that $g$ is divisible by $f^{ \alpha_0 m' }$ in $S$ and $\ord_{E_s} (g) = N$. 
\end{lem}

\begin{proof}
    Fix $m'$ divisible enough so that $m' \alpha_0$ is an integer. For any $A$-algebra $B$,  consider the following $B$-submodules of $ S_{B, mm'} := S_{A, mm'} \otimes_A B$:
    \[ M_B := f_B ^{m' \alpha_0} S_A \cap S_B = \text{im} (M_A \otimes_A B \to S_B),\]
    and similarly,
    $N_B = (\cF ^{N } _{E_A} S_{A, mm'}) \otimes _A B$, and $N'_B = (\cF _{E_A} ^{> N } S_{mm'}) \otimes_A B$. Recall that since we know that the associated graded modules of the filtration $\cF_{E_A}$ are flat over $A$, we have $N_B/N'_B = (N_A/N'_A) \otimes_A B$. Let $\psi_B$ denote the following natural map of $B$-modules:
    \[ \psi_B: M_B \cap N_B \to N_B/N'_ B. \]
    Note that with this notation, the lemma is equivalent to the statement that if the image of $\psi_ \C$ is non-zero then the image of $\psi _{\kappa(s)}$ is non-zero.

    Next, suppose $A \to k \to K$ is an extension such that $k$ and $K$ are both fields. Then, note that the image of $\psi_k$ is non-zero if and only if the image of $\psi_K$ is non-zero. In fact, we have $\psi_K = \psi_k \otimes _k K$. This is because taking intersections of subspaces and images of maps of vector spaces commute with base-change over the field. Now, setting $L = \text{Frac}(A)$, we apply the above observation to the extension $A \to L \to \C$ to conclude that if the image of $\psi_\C$ is non-zero, then, the same is true for $\psi_L$ as well.
    
    Next, we blow up the regular local ring $A$ at the maximal ideal $\fm$, and let $(A', t, K)$ be the local ring at the generic point of the exceptional divisor of the blow-up. Then, note that $A'$ is a discrete valuation ring with uniformizer $t$ such that $\fm A' = (t)$. Therefore, the residue field $K$ is an extension of $k = A/\fm$, and thus is of characteristic $p$ as well. Applying the above observation to the extension $k \to K$, we see that it is enough to show that the image of $\psi_K$ is non-zero.

    
    Now, since we know that the image of $\psi_L$ is non-zero, and $L =\text{Frac}(A')$, there exists elements $g \in S_{A', mm'}$, $h \in S_{A', mm' - \alpha_0 m'}$ and an integer $u$ such that in $S_L$, we have
    \[ g = f_{A'} ^{\alpha_0 m'} \, \frac{h}{t^{u}}, \]
    and $g$ is contained in $N_L \setminus N' _L$. Furthermore, by increasing $u$ if required, we may assume that the $t$-adic valuation of $g$ is zero. Thus, we have
    \[ g = f_{A'} ^{ \alpha _0 m'} \frac{h}{t^u}. \]
    Now, since by construction $f_k ^{\alpha_0 m'}$ is non-zero and the ring $S_K$ is an integral domain (since it is normal and $\N$-graded), we have 
    \[ 0 \neq f^{\alpha_0  m'} _{K} = f^{\alpha_0  m'} _{A'} \bmod{t}, \] 
    the $t$-adic valuation of $\frac{h}{t^{u}}$ is zero. Since $h$ is homogeneous, we may assume that $u = 0$. So we get that 
    $g = g = f_{A'} ^{ \alpha _0 m'} h$
    where $g \in S_{A', mm'}$ is non-zero modulo $t$. Therefore, the image of $g$ modulo $t$ is a non-zero element of $M_K$. Furthermore, since 
    $S_{A', mm'}/N_{A'}$ is free over $A'$, and $g$ is contained in $N_L \setminus N'_L$, $g$ must be contained in $N_{A'} \setminus N' _{A'}$. Finally, the image of $g$ modulo $t$ must be contained in $N_K/N'_K$, since we know that 
    \[ N_K/N'_K = \left(N_{A'}/N'_{A'} \right) \otimes _{A'} K =\left( \left(N_{A}/N'_{A} \right) \otimes _{A} k \right) \otimes_k K.  \]
    Since the image of $g$ modulo $t$ was already contained in $M_K$ by construction, we have that the image of $\psi_K$ is non-zero. This completes the proof of the lemma. 
\end{proof}

\end{document}
